% =====================================================================
%  Results — Characterisation of the reverse-transcriptase mining corpus
%  exports_v2, 2026-09-17.
%
%  Built from current_thesis_section_content.md. All prose and argument
%  retained; every number re-verified against the workbench tables and
%  corrected where it disagreed. See exports_v2/CORRECTIONS_current_thesis_section.md
%  for the line-by-line justification of each change.
%
%  Graphics path convention kept as in your draft: figures_db_thesis_section/*.png
%  Figure files to copy there are listed in exports_v2/FIGURE_SET.md.
%
%  Populations are named by the ids of the registry (Table~\ref{tab:units}),
%  which is computed in the workbench notebook, section 0.2.
% =====================================================================

\section{Characterisation of the reverse-transcriptase mining corpus}
\label{sec:dbchar}

The discovery pipeline described in Section~\ref{sec:discovery-pipeline} produced \num{3059700}
raw RT-system records. A record represents a reported candidate system, not an experimentally
established retron. The following analyses distinguish repeated observations, sequence diversity
and association evidence to define resources for downstream work.

Because the same biological entity can be counted in several different units, and because the
conversion between those units is not a constant, every result below names the population it was
measured on. Those populations are defined once in Table~\ref{tab:units} and referred to
thereafter by their identifiers; each figure repeats its own population and size along its lower
edge.

\subsection{What constitutes one observation}
\label{sec:dbchar:units}

The corpus cannot be counted in a single unit, because the same biological entity can appear
repeatedly for several distinct reasons. Distinguishing these forms of repetition is therefore a
prerequisite for interpreting every subsequent number. A \emph{locus} is the interval containing an
RT on a specified contig and strand, represented by the contig accession together with the start
and end coordinates and the strand. Two deposits of the same physical position under different
contig accessions are consequently treated as two accession-defined loci. A \emph{physical locus}
instead collapses supported alternative accession representations of the same contig, whereas a
\emph{distinct RT protein} is defined by its exact amino-acid sequence
(Figure~\ref{fig:units}, Table~\ref{tab:units}).

\begin{figure}[htbp]
\centering
\includegraphics[width=0.85\textwidth]{figures_db_thesis_section/A1_unit_ladder.png}
\caption[The unit ladder]{The corpus counted in each of its units. The reduction from raw records
to distinct proteins is not a quality filter: it is a change of question. Each rung answers a
different one, and no rung is the correct one in general.}
\label{fig:units}
\end{figure}

\paragraph{Redundancy of deposition.} This arises because the same assembly can be represented in
more than one collection. The corpus contains \num{111394} normalised genome identifiers
represented in two collections, while \num{203921} of the \num{2847312} accession-defined loci
(\textsf{LOC}) are represented in multiple collections. These are identifier- and locus-level
overlaps respectively, and should not be interpreted as counts of duplicate whole genomes. At the
locus level, GTDB re-publishes assemblies obtained from NCBI, producing substantial overlap between
the two bacterial collections. Consolidating supported alternative representations reduces the
\num{2847312} accession-defined loci to \num{2475684} physical loci. This reduction comprises
\num{371628} matched RefSeq/GenBank locus pairs for which the local DNA sequence, RT protein
sequence, coordinates and strand are identical. Matching RT-containing regions supports
consolidation of the locus representations, but does not establish that the corresponding whole
genomes are identical; such a claim would require assembly-level comparison.

\paragraph{Redundancy of sampling.} This has a much larger effect and is not removed by
accession-level correction. Three species account for \SI{47}{\percent} of all NCBI-assigned
records in the sequence population: \emph{Escherichia coli} (\num{449992} records),
\emph{Salmonella enterica} (\num{374290}) and \emph{Klebsiella pneumoniae} (\num{237479}). The
effect becomes apparent when organisms are compared using distinct protein sequences rather than
records. \emph{E.\ coli} contributes \num{27998} distinct RT proteins from \num{449992} records,
approximately six distinct proteins per hundred records, whereas \emph{uncultured Lachnospiraceae}
contributes \num{7539} distinct proteins from \num{16408} records, or approximately forty-six per
hundred (Figure~\ref{fig:taxa}). A taxon that is common at the record level may therefore be much
less common at the level of distinct proteins, simply because the organisms carrying it have been
sequenced more often. The same effect separates the two domains: bacterial records
carry \num{6.29} records per distinct protein against \num{1.63} for archaeal ones, and in five of
the six families with appreciable archaeal representation the archaeal ratio is the lower of the
two (Figure~\ref{fig:domain}). Archaea are only \SI{0.34}{\percent} of the corpus, so this is a
contrast rather than an archaeal survey; it is shown because it separates sequencing effort from
biology more cleanly than any within-bacterial comparison can.

\begin{figure}[htbp]
\centering
\includegraphics[width=0.85\textwidth]{figures_db_thesis_section/N6a_taxa_redundancy.png}
\caption[Sampling redundancy across species]{The most frequently deposited species carry few
distinct RT proteins; labels give distinct proteins per hundred records. Counted on the
sequence-level population \textsf{RT-SEQ}. Species-level counts use the NCBI taxonomy block only;
NCBI and GTDB ranks are not interchangeable and are never pooled.}
\label{fig:taxa}
\end{figure}
\begin{figure}[htbp]
\centering
\includegraphics[width=\textwidth]{figures_db_thesis_section/Z2_redundancy_by_domain.png}
\caption[Redundancy of deposition by taxonomic domain]{Redundancy of deposition is a property of
how much an organism has been sequenced, not of the reverse transcriptases it carries. Left:
records and distinct proteins per domain, with the ratio between them; archaeal records are close
to one distinct protein each. Right: the same families in both domains. The archaeal ratio is the
lower of the two in five of six families; \texttt{RVT-CRISPR} inverts it, on 73 archaeal proteins.
Archaeal sample sizes are printed under each family because they are small.}
\label{fig:domain}
\end{figure}

\paragraph{Consequence.} These distinctions determine what each summary statistic describes.
Record counts primarily describe deposition and sampling intensity, whereas counts of distinct
proteins describe sequence-level diversity; neither is intrinsically the more correct unit, and
each result below names the unit it uses. Exact amino-acid sequence deduplication over the whole
corpus yields \num{501561} RT proteins (\textsf{RT-BASE}). Here, exactness means equality of the
stored amino-acid sequences: a single changed residue, an ambiguity character, or a length
difference resulting from a truncated gene call is sufficient to produce a separate entry.
Exact-sequence counts are consequently \textbf{upper bounds} on biological distinctness, and
similarity-based clustering is required before interpreting any diversity estimate or partitioning
the data for downstream modelling.

\paragraph{What the exact-sequence grain measures.} How loose the upper bound is can be measured
rather than asserted, and the answer determines which claims are available. Of the \num{78287}
proteins of the Retron family (Section~\ref{sec:dbchar:families}), \num{46887}
(\SI{59.9}{\percent}) occur in exactly one genome. That
proportion is not an artefact of the two obvious technical causes: removing proteins seen only in
metagenome-assembled genomes leaves \SI{56.4}{\percent}, and additionally removing every protein
whose calls are all partial leaves \SI{56.6}{\percent} (Table~\ref{tab:singletons}). The sequences
are real.

They are nonetheless not a measure of diversity. \num{11072} of these once-seen proteins occur in
species for which at least a thousand genomes have been sequenced, and a sequence observed once
among a thousand assemblies of one species is a strain-level variant whatever its exact-sequence
status. The point is clearest within a single organism: \emph{Escherichia coli} contributes
\num{6628} distinct exact Retron RT sequences across \num{82695} genomes, and
\emph{Salmonella enterica} \num{5187} across \num{204502}, against a literature describing
retron types in these organisms in the order of ten.%
% CITATION NEEDED: a retron-classification review supporting "order of ten types" in E. coli.
 The exact-sequence grain is therefore two to three
orders of magnitude finer than biological type.

Two consequences follow, and both are load-bearing for the rest of this work. Exact-sequence
counts may be reported as counts of distinct sequences and never as counts of distinct retrons.
And no saturation, rarefaction or richness estimate may be computed at this grain. The frequency
spectrum invites one --- \num{46887} sequences are seen in a single genome and \num{17603} in
exactly two --- and a Chao1 estimator built on those two counts returns \num{140731}, some
\num{1.8} times the observed total. That figure is reported here only to be set aside: at
exact-sequence resolution it measures single-residue variation between strains of the same
species, not unsampled biology. Both claims become available after clustering, and neither
before it.

\begin{table}[htbp]
\centering\small
\caption[Proteins observed in a single genome]{Proteins observed in exactly one genome, and what
survives removing the technical explanations. The proportion is stable, so these are real
sequences; the lower block shows that they are nonetheless concentrated in the most heavily
sequenced species, where a once-seen sequence is a strain variant.}
\label{tab:singletons}
\begin{tabular}{@{}lrrr@{}}
\toprule
Stratum & Proteins & Seen once & \% \\
\midrule
All Retron proteins                          & \num{78287} & \num{46887} & 59.9 \\
\multicolumn{4}{@{}l}{\quad\emph{of which seen in exactly two genomes: \num{17603}}} \\
\quad excluding metagenome-only proteins     & \num{69240} & \num{39053} & 56.4 \\
\quad and excluding all-partial calls        & \num{45271} & \num{25642} & 56.6 \\
\midrule
\multicolumn{4}{@{}l}{\emph{Where the once-seen proteins live} (NCBI-assigned; by how well their species is sampled)} \\
\quad species with $\ge$\,1000 genomes sequenced & & \num{11072} & \\
\quad species with 100--999                      & & \num{11091} & \\
\quad species with 10--99                        & & \num{9677}  & \\
\quad species with $<$\,10                       & & \num{5798}  & \\
\bottomrule
\end{tabular}
\end{table}


\begin{table}[htbp]
\centering\small
\caption[Analytical populations]{The analytical populations used throughout this chapter. Rows are
counting units, not quality filters. Sizes are recomputed from the primary data rather than
recorded by hand.}
\label{tab:units}
\begin{tabularx}{\linewidth}{@{}llXr@{}}
\toprule
ID & Unit & Definition & Count \\
\midrule
\textsf{REC-ALL}    & raw record        & candidate-system observation as mined          & \num{3059700} \\
\textsf{REC-DIST}   & distinct record   & identical repeated records consolidated       & \num{3051238} \\
\textsf{LOC}        & locus             & RT interval on a named contig and strand      & \num{2847312} \\
\textsf{PLOC}       & physical locus    & supported duplicate representations consolidated & \num{2475684} \\
\textsf{GEN}        & genome identifier & normalised identifier in the corpus           & \num{1542433} \\
\midrule
\textsf{RT-BASE}    & exact RT protein  & distinct stored amino-acid sequence, whole corpus & \num{501561} \\
\textsf{RT-BASE-1F} & exact RT protein  & \quad of which carrying one family label      & \num{493956} \\
\textsf{RT-MULTI}   & exact RT protein  & \quad of which multi-labelled                 & \num{7593} \\
\textsf{RT-CROSS}   & exact RT protein  & \quad of which cross-labelled                 & 12 \\
\textsf{RT-SEQ}     & exact RT protein  & sequence-level population (Sec.~\ref{sec:dbchar:qc}) & \num{493964} \\
\textsf{RT-CTX}     & exact RT protein  & context-level population                      & \num{477956} \\
\midrule
\textsf{PL-ALL}     & placement         & one RNA call at one RT locus                  & \num{346722} \\
\textsf{PL-CANON}   & placement         & eligible and non-redundant                    & \num{344154} \\
\textsf{NC-ALL}     & exact RNA         & distinct called RNA sequence                  & \num{16466} \\
\textsf{PAIR-ELIG}  & exact pair        & distinct (RT, RNA) over eligible placements   & \num{30924} \\
\textsf{PAIR-CANON} & exact pair        & distinct (RT, RNA) over \textsf{PL-CANON}     & \num{30427} \\
\bottomrule
\end{tabularx}
\end{table}

\subsection{Annotation artefacts and what survives them}
\label{sec:dbchar:qc}

\num{31504} records carry no annotated coding sequence for the RT they are anchored on. Resolving
them separates a recoverable annotation failure from an irrecoverable one
(Table~\ref{tab:recovery}). For \num{14188} records the coding sequence was reconstructed by
translating the genomic window and matching the stored protein, and these retain full genomic
context. For \num{16688} the protein is usable but the genomic context is not: the RT lies wholly
beyond the retrieved contig (\num{9128}), crosses a contig boundary (\num{3636}), or sits in an
inverted window from which no sequence could be recovered (\num{3924}). A further 628 are
irrecoverable.

\begin{table}[htbp]
\centering
\caption[Records lacking an annotated RT coding sequence]{Resolution of the \num{31504} records
lacking an annotated RT coding sequence. The context-less classes derive entirely from
metagenome-assembled genomes.}
\label{tab:recovery}
\begin{tabular}{llrl}
\toprule
Outcome & Basis & Records & Enters \\
\midrule
Coding sequence recovered   & reconstructed and verified against the window & \num{14188} & both populations \\
Protein only, no context    & RT beyond the retrieved contig                & \num{9128}  & \textsf{RT-SEQ} only \\
                            & RT crossing a contig boundary                 & \num{3636}  & \textsf{RT-SEQ} only \\
                            & inverted window, no sequence retrieved        & \num{3924}  & \textsf{RT-SEQ} only \\
Irrecoverable               & geometry ill-posed                            & 628         & neither \\
\bottomrule
\end{tabular}
\end{table}

This is the reason the analysis carries two populations rather than one. A record whose protein is
sound but whose neighbourhood is unknown is evidence about sequence and no evidence about
architecture. It therefore enters the sequence-level population \textsf{RT-SEQ}
(\num{493964} distinct proteins) and is excluded from the context-level population \textsf{RT-CTX}
(\num{477956}); the \num{16008} proteins that separate the two are usable sequences without
defensible genomic context, and discarding them would have removed real diversity to buy a
geometry guarantee that only geometry needs.

All \num{16688} context-less records originate from metagenome-derived collections and none from
isolate-derived collections, where the rate is zero; within the human-gut MAG collection they are
\SI{10.81}{\percent} of records. Short, fragmented assemblies simply do not extend far enough for a
gene near a contig edge to be fully represented. This is an assembly property rather than a
detection failure.

Windows truncated at a contig edge were retained rather than excluded, because they are numerous
(approximately \SI{41}{\percent} of records carry a truncation flag) and their RT coordinates are
verified; where truncation affects a measurement it is handled by stratification, as in
Section~\ref{sec:dbchar:architecture}.

\subsection{Composition of the RT family landscape}
\label{sec:dbchar:families}

The exact-sequence baseline over the whole corpus contains \num{501561} distinct RT proteins
(\textsf{RT-BASE}). These divide into \num{493956} proteins carrying a single family label
(\textsf{RT-BASE-1F}), \num{7593} multi-labelled proteins (\textsf{RT-MULTI}), and 12 proteins that
appear under more than one single-family label in different records (\textsf{RT-CROSS}). The
record-filtered sequence population \textsf{RT-SEQ} is a different selection of the same proteins:
it applies the record-level eligibility rules of Section~\ref{sec:dbchar:qc} and, by construction,
\textbf{excludes the multi-labelled stratum entirely}. It comprises \num{493952} single-label
proteins and the 12 cross-labelled ones, \num{493964} in total, across 41 family labels. The two
counts must not be differenced: \textsf{RT-SEQ} is not \textsf{RT-BASE} minus \textsf{RT-MULTI}.

Three families dominate protein-level diversity. Of the \num{493956} single-family baseline
proteins, RVT-GII accounts for \SI{51.95}{\percent}, Retron for \SI{15.85}{\percent} and RVT-DGRs
for \SI{15.41}{\percent}, together approximately \SI{83}{\percent}. The Retron subset contains
\num{78287} exact proteins.

Reporting each family on records and on distinct proteins side by side makes the redundancy of
Section~\ref{sec:dbchar:units} concrete and family-specific (Figure~\ref{fig:families}). Families
whose record share greatly exceeds their protein share are those concentrated in repeatedly
deposited isolate genomes; the ratio of records to distinct proteins, shown alongside, is a direct
measure of how often the same protein was deposited. It ranges from 1.8 to 58.6 across families,
which is itself the reason a single corpus-wide redundancy factor would be misleading.

\begin{figure}[htbp]
\centering
\includegraphics[width=\textwidth]{figures_db_thesis_section/N2_family_composition_all.png}
\caption[RT family composition]{All 41 single-family labels, shown as a share of records and of
distinct proteins. The gap between the two bars is the redundancy of deposition; the right panel
gives it directly as records per distinct protein.}
\label{fig:families}
\end{figure}

\paragraph{Length distribution.} RT protein length differs between families but does not separate
them, and the comparison is only meaningful once call completeness is accounted for. Completeness
is derived from the gene caller's partial-gene attribute, which records separately whether the
5$'$ and 3$'$ ends of the called gene were found or ran off the contig. Across distinct records,
\SI{88.19}{\percent} of RT genes are complete at both ends, \SI{4.79}{\percent} are truncated at
the 5$'$ end, \SI{4.31}{\percent} at the 3$'$ end and \SI{1.68}{\percent} at both; the remaining
\SI{1.03}{\percent} are the records of Table~\ref{tab:recovery} that carry no RT coding sequence at
all. Since \SI{25.2}{\percent} of distinct proteins are partial at every occurrence, and partial
calls are short by construction, a length distribution pooling both classes understates every
family. Figure~\ref{fig:lengths} therefore reports complete and partial calls separately.

On complete calls the overall median is 420~aa, with Retron at 366~aa and RVT-GII at
434~aa; pooling the two completeness classes moves the same three figures down to 385, 342
and 407~aa respectively, which is the size of the artefact being corrected. Several
ungrouped families occupy substantially longer regimes --- RVT-UG5 at a median of 1003~aa
and RVT-UG8 at 665~aa --- consistent with domain fusion. Extreme values were retained rather
than trimmed.

\begin{figure}[htbp]
\centering
\includegraphics[width=\textwidth]{figures_db_thesis_section/N3_rt_length_by_completeness.png}
\caption[RT protein length by family and completeness]{RT protein length by family, separated by
call completeness, for the twelve families with the most complete calls. Bars are the
interquartile range, whiskers Tukey, gold the median; the table gives complete calls only. Partial
calls are short by construction, so pooling the two classes understates every family.}
\label{fig:lengths}
\end{figure}

\paragraph{Multi-labelled proteins.} Re-scoring the \num{7593} multi-labelled proteins against the
family profile library places the highest-scoring family among the labels already assigned in
\SI{99.91}{\percent} of cases, so the multiple labels are not erroneous. They are nonetheless not
resolvable: the median margin between the best- and second-best-scoring family is 5~bits,
and \SI{99.8}{\percent} of these proteins have a margin below 20~bits, against a median of
81.7~bits in a seeded single-family control in which the best-scoring profile matched the
known label for \SI{98.10}{\percent} of proteins. Assigning each protein to its highest-scoring
family would therefore distribute \num{7593} proteins among families on evidence that does not
distinguish those families. They are retained as their own stratum with their label sets preserved:
they enter no family denominator, because they would otherwise be counted in two families at once,
and they are reported separately wherever sequence-level diversity is at issue.

\paragraph{Cross-labelled proteins.} The 12 proteins that appear under two different single-family
labels are not scattered across the classification: \textbf{all twelve are
\texttt{RVT-CRISPR}\,/\,\texttt{RVT-CRISPR-like}} (Table~\ref{tab:crosslabel}). This is a boundary
between two adjacent labels rather than a general instability of the family assignment. The two
labels were not merged, since their length distributions differ substantially --- a pooled median
of 468 against 331~aa, and 473 against 414~aa on complete calls alone; both original
labels are preserved on each protein.

\begin{table}[htbp]
\centering\small
\caption[Cross-labelled RT proteins]{The 12 exact RT proteins observed under more than one
single-family label. Every case lies on the same label boundary.}
\label{tab:crosslabel}
\begin{tabular}{lrrrr}
\toprule
Label set & Proteins & Length range (aa) & Records & Genomes \\
\midrule
\texttt{RVT-CRISPR} \textbar{} \texttt{RVT-CRISPR-like} & 12 & 202--983 & 255 & 144 \\
\bottomrule
\end{tabular}
\end{table}

\subsection{Detection of the associated non-coding RNA}
\label{sec:dbchar:ncrna}

Non-coding RNAs were detected with a library of 21 covariance models, yielding \num{346722} calls
(\textsf{PL-ALL}) that reduce to \num{16466} distinct RNA sequences (\textsf{NC-ALL}). Each
distinct sequence is matched by exactly one model, so per-model counts partition the sequence set
cleanly (Figure~\ref{fig:ncrna}).

Two of the 21 models are named \texttt{OutgroupA} and \texttt{OutgroupB}, and their share of the
sequence set is far larger than their share of calls: \texttt{OutgroupA} alone matches \num{3319}
of the distinct RNA sequences and \texttt{OutgroupB} a further 338, together \SI{22}{\percent} of
the sequence set against only \SI{3.7}{\percent} of calls. The name is a property of the model's
place in the library, however, and not a verdict on the hits it makes. Tested against the
placements themselves, outgroup-model calls at Retron loci sit in the same architecture as the
other nineteen: \SI{93.1}{\percent} upstream against \SI{94.6}{\percent}, \SI{99.3}{\percent}
same-strand against \SI{99.1}{\percent}, \SI{99.2}{\percent} with no intervening coding sequence
against \SI{94.3}{\percent}, and a median separation of 27~bp against 55~bp --- if anything
tighter than the rest. They are therefore retained, and the per-model composition is reported so
that a downstream user may stratify on it. What the disparity does require is that any count of
distinct RNA sequences state which models it includes, because the two model groups contribute
very differently to sequences and to calls.

Over distinct sequences, detected RNAs occupy a narrow length regime (median 151~nt,
interquartile range 131--193); over calls, where heavily deposited sequences dominate, the same
statistic is 158~nt (141--167). Most individual models return hits of nearly constant length.
These are model-defined hit boundaries rather than independently determined transcript boundaries,
and the distinction matters for any downstream use of the sequences.

\begin{figure}[htbp]
\centering
\includegraphics[width=0.9\textwidth]{figures_db_thesis_section/C2_ncrna_length_by_model.png}
\caption[Detected RNA length by covariance model]{Length of detected RNAs by covariance model, on
distinct sequences. Each distinct sequence is matched by exactly one model. These are model-defined
hit boundaries, not transcript boundaries.}
\label{fig:ncrna}
\end{figure}

Hit confidence is high enough to be used as a filter rather than merely reported:
\SI{97.74}{\percent} of calls have an expectation value at or below $10^{-5}$ and
\SI{55.81}{\percent} at or below $10^{-20}$ (Figure~\ref{fig:cmconf}). Confidence is not evenly
distributed across models, however. The worst-affected model, \texttt{TypeIIIA1}, returns
\SI{52.0}{\percent} of its \num{1647} calls above $10^{-5}$, and \texttt{TypeIIIA2}
\SI{39.2}{\percent} of \num{3215}. A global threshold is therefore not a neutral operation: it
removes a few percent of the corpus overall while removing half of some models.

\begin{figure}[htbp]
\centering
\includegraphics[width=\textwidth]{figures_db_thesis_section/N6b_cm_confidence.png}
\caption[Covariance-model hit confidence]{Left: hit confidence over eligible, non-redundant
placements; a threshold at $E \le 10^{-5}$ retains \SI{97.74}{\percent} of calls. Right: the same
threshold does not cost every model equally --- the eight most affected models, worst first.}
\label{fig:cmconf}
\end{figure}

Detection is strongly family-restricted: \SI{52.77}{\percent} of Retron loci carry at least one
call, against below \SI{0.5}{\percent} for every other family (the highest being RVT-UG10 at
\SI{0.48}{\percent}). This asymmetry is expected and is not evidence about the biology of other
families. The models are retron models, so their failure to match elsewhere measures the scope of
the instrument. Figure~\ref{fig:famxmodel} accordingly reports what the retron models find across
families, and should not be read as a statement of which families lack associated RNAs.
Establishing absence would require a positive control demonstrating that the same instrument
recovers a known-present case on an appropriate substrate.

\begin{figure}[htbp]
\centering
\includegraphics[width=\textwidth]{figures_db_thesis_section/N4_family_x_model.png}
\caption[Retron covariance models across RT families]{What the retron covariance models find across
RT families. Near-absence outside the Retron family reflects the scope of the model library and is
not evidence that other families lack associated RNAs.}
\label{fig:famxmodel}
\end{figure}

\subsection{Genomic architecture of the RT--RNA association}
\label{sec:dbchar:architecture}

Geometry was computed from coordinates and strand throughout, and the shipped
\texttt{position\_relative\_to\_rt} field was not used: it is null on \num{331897} placements and
matches no coordinate frame, and is retained as provenance only. The \num{346722} placements reduce
to \num{344154} after removing placements on duplicated records and repeated observations of the
same call at one locus; these are referred to below as \emph{eligible, non-redundant placements}
(\textsf{PL-CANON}), a technical description that carries no implication of biological validity.
Every geometry figure in this section is measured on that one population.

The dominant arrangement is close, same-strand and 5$'$ (Table~\ref{tab:geometry}): the RNA lies
upstream of the RT in \SI{94.51}{\percent} of placements, overlaps it in \SI{3.52}{\percent} and
lies downstream in \SI{1.96}{\percent}; the median upstream separation is 55~bp;
\SI{99.12}{\percent} are on the same strand as the RT; and \SI{94.41}{\percent} have no annotated
coding sequence between the two features.

These properties are strongly coupled rather than independently distributed
(Figure~\ref{fig:joint}). The single combination of upstream, same strand, no intervening coding
sequence and separation within 200~bp accounts for \SI{55.22}{\percent} of placements, and
the remaining combinations fall away by orders of magnitude. Re-weighting from placements to
distinct physical loci leaves the marginals almost unchanged --- upstream moves from
\SI{94.51}{\percent} to \SI{94.91}{\percent} --- so the arrangement is not produced by one position
being deposited repeatedly. Re-weighting to distinct sequence pairs is more discriminating: the
band within 200~bp strengthens to \SI{73.86}{\percent} of distinct pairs, while the
1--5~kb band collapses from \SI{32.83}{\percent} of placements to \SI{2.95}{\percent}
of pairs, showing the latter to be carried by few sequence pairs observed many times. Close 5$'$
adjacency is therefore the dominant arrangement under every weighting, and the only band that
strengthens when redundancy is removed.

\begin{figure}[htbp]
\centering
\includegraphics[width=\textwidth]{figures_db_thesis_section/D6_joint_geometry.png}
\caption[Joint geometry and the effect of weighting]{Direction, strand, adjacency and separation
are strongly coupled. The same combinations re-weighted from placements to distinct physical loci
and to distinct sequence pairs: only the close upstream band strengthens when redundancy is
removed, which is what distinguishes it from a deposition artefact.}
\label{fig:joint}
\end{figure}

These are coordinates. They are compatible with a shared transcriptional unit and do not
demonstrate one, and they carry no information about recognition between a particular RT and a
particular RNA.

\paragraph{Overlapping placements.} The \num{12124} overlapping placements are a benign class
rather than a failure. In \SI{97.36}{\percent} of them the RNA overlaps the RT gene itself and not
some other feature; only 45 overlap a different coding sequence. The extent is small: a median of
22~bp, with \SI{58.5}{\percent} overlapping a tenth or less of the RNA, and 266 placements
in which the RNA lies fully enclosed within a coding sequence. This is the short encroachment
expected where an RNA abuts the 5$'$ end of a gene in a compact locus. Overlapping placements are
accordingly retained in the association resource as a flagged class, and only the small
pathological subsets are set aside.

\paragraph{A technical mode among downstream placements.} Where the extraction window begins at the
start of a contig there is no sequence upstream of the RT to search, so any RNA found is
necessarily reported downstream. These forced placements accumulate at one separation: \num{5234}
placements, \SI{77.4}{\percent} of all downstream placements, lie within 150~bp of a median
of 2682~bp, involve only 18 distinct RNA sequences across 15 species, and are
\SI{99.81}{\percent} truncated at the contig start. Stratifying on that truncation rather than on
the mode isolates the population the artefact does not explain: the \num{1041} downstream
placements that are \emph{not} contig-start-truncated lie at a median of 64~bp over 149
distinct sequences --- as close to the RT as upstream placements are, and in the opposite
orientation.

\paragraph{Placements on the opposite strand.} The \num{3041} placements on the strand opposite the
RT are concentrated rather than scattered: \SI{84.5}{\percent} come from one covariance model
(\texttt{TypeXIIIA\_firmi}) and \SI{96}{\percent} occur at Retron loci. They differ from
same-strand placements in every respect that matters. Their median separation is 4750~bp
against 44~bp, and only \SI{12.5}{\percent} have no intervening coding sequence against
\SI{95.1}{\percent}. They are not an artefact of repeated deposition, spanning \num{1324} physical
loci across 67 species, and they are not weak matches, since \num{2921} of \num{3041} have an
expectation value at or below $10^{-5}$. The same model additionally produces \num{1576}
same-strand placements at a median of 64~bp across 292 species, so the contrast is internal
to one model and not a property of the model library. A hit several kilobases away, with genes in
between, is a match elsewhere in the retained genomic window rather than a component of the RT's
own locus, and these placements are excluded from the association resource on that basis.

\paragraph{Retron models at non-Retron RT loci.} Non-Retron RT loci carrying a retron covariance
model hit show both a weaker positional preference (\SI{57.42}{\percent} of the 256 such
placements are upstream, against \SI{94.54}{\percent} at Retron loci) and a far greater separation:
a median of 1802~bp for upstream placements and 4281~bp for downstream ones, against
55~bp at Retron loci (Figure~\ref{fig:architecture}). They are not architecturally consistent
with retron loci and are retained as candidates for annotation follow-up rather than described as
divergent retrons.

\begin{figure}[htbp]
\centering
\includegraphics[width=\textwidth]{figures_db_thesis_section/N5_architecture_retron_vs_not.png}
\caption[RT--RNA architecture]{Architecture of the RT--RNA association, comparing Retron loci with
non-Retron loci carrying a retron model hit, on eligible non-redundant placements. Overlapping
placements are shown as their own category; the rightmost panel gives the extent of overlap, which
is predominantly a short encroachment on the RT gene itself.}
\label{fig:architecture}
\end{figure}

\begin{table}[htbp]
\centering
\caption[Geometry of eligible non-redundant placements]{Geometry of the \num{344154} eligible,
non-redundant RT--RNA placements (\textsf{PL-CANON}).}
\label{tab:geometry}
\begin{tabular}{lrr}
\toprule
Property & Placements & Percentage \\
\midrule
RNA upstream of the RT   & \num{325269} & 94.51 \\
Overlapping the RT       & \num{12124}  & 3.52 \\
RNA downstream of the RT & \num{6761}   & 1.96 \\
\midrule
Same strand as the RT    & \num{341113} & 99.12 \\
No intervening CDS       & \num{324904} & 94.41 \\
One intervening CDS      & \num{12427}  & 3.61 \\
\midrule
\multicolumn{3}{l}{Median upstream separation: 55~bp} \\
\bottomrule
\end{tabular}
\end{table}

\subsection{Structure of the association at sequence level}
\label{sec:dbchar:pairs}

Reducing placements to distinct pairs of sequences requires choosing which placement population to
reduce, and the two available choices do not give the same answer (Table~\ref{tab:pairbasis}). The
registered pair view is built on eligible placements and contains \num{30924} pairs
(\textsf{PAIR-ELIG}); only \num{30427} of those are derivable from the eligible, non-redundant
placements used for geometry (\textsf{PAIR-CANON}). The \num{497} pairs that separate them are
represented solely by placements removed during de-duplication. Pair topology is reported below on
the registered \num{30924}, and the association resource of Section~\ref{sec:dbchar:dataset} is
built on the \num{30427}; the two are never mixed.

\begin{table}[htbp]
\centering\small
\caption[Basis of the exact-pair view]{Distinct RT--RNA pairs depend on which placement population
they are reduced from. The registered view matches the eligible population.}
\label{tab:pairbasis}
\begin{tabular}{@{}lrrr@{}}
\toprule
Placement population & Pairs & Distinct RTs & Distinct RNAs \\
\midrule
All placements                         & \num{31029} & \num{29297} & \num{16466} \\
Eligible (\textsf{PAIR-ELIG})          & \num{30924} & \num{29192} & \num{16458} \\
Eligible, non-redundant (\textsf{PAIR-CANON}) & \num{30427} & \num{28976} & \num{16011} \\
\bottomrule
\end{tabular}
\end{table}

On the registered view --- \num{30924} pairs over \num{29192} distinct RT proteins and \num{16458}
distinct RNAs --- the association is, as observed, close to one-to-one: \SI{96.85}{\percent} of the
RT proteins occur with exactly one RNA sequence and \SI{82.28}{\percent} of RNAs with exactly one
RT protein. Of the \num{14918} connected components of the bipartite graph, \num{12079}
(\SI{81}{\percent}) are strictly one-to-one.

The exceptions are, on inspection, consequences of exact-sequence identity rather than of biology,
and they illustrate why the caveat of Section~\ref{sec:dbchar:units} is load-bearing. A pair here
is a distinct combination of one exact RT sequence and one exact RNA sequence; because both sides
are exact, a single trimmed nucleotide produces a second pair. Where an RT protein has several RNA
partners, those partners are almost always length variants of one RNA: for 659 of the 921 such
proteins the partner sequences differ in length by 5~nt or less, and the most extreme case,
176 partners, is a near-continuous length series from 68 to 223~nt, produced by the model
trimming its hit slightly differently at each locus. Correspondingly, the RNA with the greatest
number of RT partners (705) occurs across \num{25748} genomes but only 17 species, and its partners
range from 58 to \num{1254} amino acids within one family --- that is, length and truncation
variants of related proteins in heavily sequenced taxa.

Apparent many-to-many structure is therefore largely redundancy, and the partner counts above are
upper bounds. Establishing the true topology requires clustering both sequence sets at a stated
identity threshold, or trimming RNAs to a common model-defined boundary, before partners are
counted. That clustering has not been performed.

\paragraph{Recurrence.} Repetition was classified by the kind of repetition it represents, which
separates re-deposition from independent observation. Of the \num{30924} pairs, \num{12005} occur
once; \num{9362} recur only as multiple database copies of a single physical locus; \num{4452}
recur across genomes of one species; \num{5056} recur across species; and 49 fall into two small
residual classes. Roughly one third of the apparent recurrence in this resource is therefore
deposition rather than independent observation. Only the last two classes --- \num{9508} pairs ---
constitute repeated observation of an association in distinct genomic backgrounds, and only those
support an evolutionary reading. A pair seen a hundred times in one re-deposited assembly is one
observation, not a hundred.

\subsection{From observations to an association resource}
\label{sec:dbchar:dataset}

The preceding sections each identify a population that should not enter a resource intended to
represent RT--RNA association, and applying them in sequence gives the reduction shown in
Figure~\ref{fig:funnel} and Table~\ref{tab:funnel}. Beginning from \num{30427} distinct sequence
pairs supported by eligible, non-redundant placements, restriction to the Retron family
(\num{30287}), to same-strand placements (\num{29837}), to placements that are not downstream
(\num{29440}), to those with no intervening coding sequence (\num{28442}), to a separation within
200~bp (\num{25673}) and to high-confidence model hits (\num{23680}) retains
\SI{77.8}{\percent} of the starting set.

\begin{figure}[htbp]
\centering
\includegraphics[width=\textwidth]{figures_db_thesis_section/N8_dataset_funnel.png}
\caption[Construction of the association resource]{Reduction from eligible, non-redundant
placements to the association resource, with the count removed at each step and the finding that
motivates it.}
\label{fig:funnel}
\end{figure}

Each reduction is attributable. Restriction to Retron follows from detection being
retron-model-scoped (Section~\ref{sec:dbchar:ncrna}); the exclusion of opposite-strand and
downstream placements follows from their being, respectively, distant single-model matches
elsewhere in the window and a window-truncation artefact
(Section~\ref{sec:dbchar:architecture}); the adjacency and separation criteria follow from the
joint geometry; and the confidence threshold costs little overall because \SI{97.74}{\percent} of
calls already satisfy it, though it falls unevenly across models. Overlapping placements are
retained, for the reasons given above.

The 200~bp separation criterion is the joint-geometry mode, not a completeness bound: it is
the single most expensive step in the funnel, removing \num{2769} pairs, and relaxing it to
500~bp recovers a substantial part of them. It is offered as the default operating point
rather than as a biological boundary, and the resource is published in tiers so that a downstream
application may select its own.

\begin{table}[htbp]
\centering
\caption[Construction of the association resource]{Construction of the RT--RNA association
resource on \textsf{PAIR-CANON}, with the finding motivating each reduction.}
\label{tab:funnel}
\begin{tabular}{llrr}
\toprule
Criterion & Motivating finding & Pairs & Retained (\%) \\
\midrule
Eligible, non-redundant   & de-duplication, coordinate verification & \num{30427} & 100.0 \\
Retron family             & detection is retron-model-scoped        & \num{30287} & 99.5 \\
Same strand               & opposite strand matches elsewhere       & \num{29837} & 98.1 \\
Not downstream            & window-truncation artefact              & \num{29440} & 96.8 \\
No intervening CDS        & operonic architecture                   & \num{28442} & 93.5 \\
Within 200~bp       & the joint-geometry mode                 & \num{25673} & 84.4 \\
Model hit $E \le 10^{-5}$ & detection confidence                    & \num{23680} & 77.8 \\
\bottomrule
\end{tabular}
\end{table}

Representative loci drawn from this resource are shown in Figure~\ref{fig:operons}.

\begin{figure}[htbp]
\centering
\includegraphics[width=\textwidth]{figures_db_thesis_section/M2_operons_strand_normalised.png}
\caption[Representative loci]{Representative RT--RNA loci selected by the criteria of this section.
All are drawn on the plus strand so that transcription runs left to right; the strand each locus
was actually found on is marked. Grey arrows are flanking genes outside the delimited run.}
\label{fig:operons}
\end{figure}

Two limitations remain, and both are deliberate. Sequence clustering has not been applied, so the
pair counts are upper bounds and the resource is not yet partitioned for modelling: variants of a
single RNA, and truncated variants of a single protein, would otherwise fall on both sides of any
train/test split. And the resource is conditioned on the detection route, as
Section~\ref{sec:dbchar:tools} shows.

\subsection{Dependence on the annotation route}
\label{sec:dbchar:tools}

Every result above is conditioned on which tools annotated a locus, and the size of that effect is
only visible on the protein axis. Across the \num{78287} distinct Retron RT proteins, the three
detection routes partition as shown in Figure~\ref{fig:venn}. myRT alone accounts for \num{29436}
distinct proteins, more than the \num{20203} recovered by all three tools together; on records the
ordering reverses (\num{126406} against \num{352705}), because the three-tool intersection is
dominated by repeatedly deposited organisms. A tool-defined subset of this corpus is not a random
subset of it, and the two axes give opposite answers to the question of which route contributes
most.

\begin{figure}[htbp]
\centering
\includegraphics[width=\textwidth]{figures_db_thesis_section/N1_tool_venn.png}
\caption[Detection route]{Agreement between detection routes on distinct Retron RT proteins (left)
and on records (centre), and the RNA carriage rate of each combination (right). A protein is
assigned to a tool set by aggregating across all of its records, so the protein panel is a true
partition; the protein and record axes give opposite orderings.}
\label{fig:venn}
\end{figure}
% GENERATED by notebook section Z2.1 - do not edit by hand.
\begin{table}[htbp]
\centering\footnotesize
\setlength{\tabcolsep}{5pt}
\caption[Detection-tool agreement on retron reverse transcriptases]{%
\textbf{Agreement between the three detection tools, and its relationship to non-coding RNA
co-occurrence.} Regions are exclusive and partition each axis. A protein is assigned to a region
by aggregating the tools across all of its records, so the two axes have different shapes: myRT
alone is the largest region by protein and the second smallest by record. RNA co-occurrence varies
4.7-fold across regions on the protein axis. Every RNA call in this
corpus was made by Infernal against the covariance-model library, for each record alike; none
originated from a detection tool, so the association is a property of what each tool's retron
model requires in order to fire, not evidence that the tool found the RNA.}
\label{tab:venn_agreement}
\begin{tabular}{@{}lrrrrrr@{}}
\toprule
& \multicolumn{3}{c}{\textbf{Per distinct record}} & \multicolumn{3}{c}{\textbf{Per distinct RT protein}} \\
\cmidrule(lr){2-4}\cmidrule(l){5-7}
\textbf{Detected by} & \textbf{n} & \textbf{\%} & \textbf{\% RNA} & \textbf{n} & \textbf{\%} & \textbf{\% RNA} \\
\midrule
myRT $+$ PADLOC $+$ DefenseFinder & \num{352705} & 53.17 & 70.39 & \num{20203} & 25.81 & 53.26 \\
myRT only & \num{126406} & 19.06 & 13.50 & \num{29436} & 37.60 & 18.64 \\
myRT $+$ DefenseFinder & \num{74697} & 11.26 & 16.97 & \num{13199} & 16.86 & 24.45 \\
myRT $+$ PADLOC & \num{61541} & 9.28 & \textbf{89.23} & \num{8994} & 11.49 & \textbf{87.58} \\
PADLOC only & \num{26227} & 3.95 & 18.63 & \num{3492} & 4.46 & 25.09 \\
PADLOC $+$ DefenseFinder & \num{13003} & 1.96 & 44.23 & \num{1232} & 1.57 & 48.94 \\
DefenseFinder only & \num{8729} & 1.32 & 14.66 & \num{1731} & 2.21 & 18.83 \\
\midrule
\textbf{Total} & \textbf{\num{663308}} & \textbf{100} & 51.99 & \textbf{\num{78287}} & \textbf{100} & 37.24 \\
\bottomrule
\end{tabular}
\end{table}

The apparent rate of RNA association depends on the same choice. Among proteins detected by myRT
and PADLOC the RNA carriage rate is \SI{87.58}{\percent}, and among those detected by myRT alone it
is \SI{18.64}{\percent}. Table~\ref{tab:venn_agreement} gives every region on both axes.
The RNA calls in this corpus were all produced by Infernal against the covariance-model library,
for every record alike; not one originated from a detection tool. The association between tool
region and RNA carriage is therefore a selection effect --- a tool whose retron model requires the
RNA-adjacent architecture only fires where such an architecture exists, and Infernal then finds the
RNA there --- and not evidence that the tool detected the RNA. The covariance models are not evenly distributed across tool combinations
either, with individual models occurring almost entirely within particular combinations, so which
model matched and which tools called the locus are not independent.

Two limitations follow. The tools share model lineage, so their agreement is not independent
corroboration; and the absence of a tool from a record cannot be distinguished from that tool never
having been run on the corresponding genome, so these figures bound tool coverage from below.

\subsection{The resulting datasets}
\label{sec:dbchar:resources}

The funnel of Section~\ref{sec:dbchar:dataset} is not a single cut-off but a nested series, and the
appropriate stopping point depends on the use. Table~\ref{tab:tiers} gives the four tiers that
result, each reported both as exact sequence pairs and as the number of RT-to-RNA-\emph{type}
associations those pairs collapse to.

\begin{table}[htbp]
\centering\small
\caption[The resulting association datasets]{The four resulting tiers. Tiers are nested:
T4 $\subset$ T3 $\subset$ T2 $\subset$ T1. ``Pairs'' counts distinct combinations of one exact RT
sequence and one exact RNA sequence; ``associations'' counts distinct combinations of one exact RT
and one covariance model, which is invariant to boundary trimming. Their difference is the
boundary-jitter inflation.}
\label{tab:tiers}
\begin{tabularx}{\linewidth}{@{}lXrrrr@{}}
\toprule
Tier & Rule added & Pairs & Assoc. & RTs & RNAs \\
\midrule
T1 observed        & eligible, non-redundant, Retron            & \num{30287} & \num{28845} & \num{28838} & \num{15906} \\
T2 architecture    & + same strand, not downstream, no intervening CDS, $\le$\,200~bp & \num{25673} & \num{24677} & \num{24677} & \num{13258} \\
T3 high-confidence & + model hit $E \le 10^{-5}$                & \num{23680} & \num{22727} & \num{22727} & \num{11849} \\
T4 independently recurrent & + recurs across genomes or species & \num{7476} & --- & \num{7187} & \num{3829} \\
\bottomrule
\end{tabularx}
\end{table}

\paragraph{Why these are the final sets.} T3 is the default. Every criterion applied to reach it is
attributable to a finding of this chapter rather than to a preference, and each was shown to remove
a population that is explicable as detector scope, window truncation or a match elsewhere in the
retained window. T1 is the permissive alternative for any analysis that must not condition on
architecture --- notably, any test of whether architecture predicts something else, which would
otherwise be circular. T4 is the subset for evolutionary questions: it retains only pairs observed
in more than one genomic background, and is therefore the only tier in which a repeated observation
is a repeated \emph{event} rather than a repeated deposition.

\paragraph{Is the many-to-many structure resolved?} Partly, and the two directions must be answered
separately. Within T3, \SI{97.35}{\percent} of RT proteins occur with exactly one RNA sequence.
Of the 602 that do not, \textbf{every one draws all of its partners from a single covariance
model}, and for 530 of them the partner sequences differ in length by 5~nt or less (563 by 20~nt or
less). These are the same RNA with a differently trimmed boundary, not competing partners.
Collapsing each RT to one association per model therefore resolves this direction completely, and
costs \num{953} of the \num{23680} pairs (\SI{4.0}{\percent}): the \num{22727} associations of
Table~\ref{tab:tiers} are exactly the \num{22727} distinct RT proteins.

The opposite direction is \textbf{not} resolved. \SI{81.18}{\percent} of RNA sequences occur with
exactly one RT protein, but the most connected RNA has 690 distinct RT partners. Those partners are
not competing RNAs seen from the other side: they span 58 to \num{1254} amino acids within a single
family label set, that is, length and truncation variants of related proteins in heavily sequenced
taxa. Resolving this requires clustering the \emph{proteins}, not the RNAs, and until that is done
every RNA-side partner count in this chapter is an upper bound.

\paragraph{Considerations before use.} Four, in the order they will bite.

\begin{enumerate}
\item \emph{Exact-sequence identity inflates both sides.} Use the association count, not the pair
count, wherever a single trimmed nucleotide should not create a new entity. The two differ by
\SI{4.0}{\percent} in T3, and by more in T1, where the architecture filters have not yet removed
the noisiest placements.

\item \emph{Clustering will reduce these numbers further, and unevenly.} On the protein side it is
immediately runnable and is the one that matters, since it is what resolves the RNA-to-many-RT
direction. On the RNA side no clustering input exists yet: the derived datasets carry
\texttt{nc\_seq\_hash} and the hit length, but not the RNA sequence itself, so exporting those
sequences is a prerequisite rather than an optional refinement. Expect the reduction to fall
hardest on the metagenome-derived entries, where near-identical variants proliferate, so cluster
counts will not be a constant fraction of exact counts and the two must never be compared across
tiers.

\item \emph{Recurrence must be counted on genomes, not on species names.} Distinct
\texttt{genome\_id\_norm} values already overstate independence, because one physical locus is
republished under several genome accessions; that is why T4 is defined on the registered recurrence
classification rather than on a recount. Counting distinct species names is worse: \num{2760} T3
pairs sit in a single genome that NCBI and GTDB name differently, and a species-based rule scores
every one of them as cross-species recurrence.

\item \emph{Every tier is conditioned on the detection route}
(Section~\ref{sec:dbchar:tools}), and the covariance models are retron models. A tier is a set of
associations that these instruments can see; it is not a census. Absence from any tier is not
evidence of absence in the genome.
\end{enumerate}

Sequence clustering therefore remains the single outstanding dependency, and it is the reason these
sets are offered as tiers rather than as one partitioned modelling dataset: variants of one RNA,
and truncated variants of one protein, would fall on both sides of any train/test split made today.

% =====================================================================
%  SUPPLEMENTARY — referenced from the text above, placed in the appendix
% =====================================================================
%
%  S1  A2_redundancy_by_database.png      redundancy ratios per source collection
%  S2  A3_occurrence_distribution.png     occurrence distribution per exact RT protein
%  S3  A4_database_attribution.png        coverage vs exclusive partition by collection
%  S4  B1_family_composition.png          family composition, baseline view
%  S5  B3_completeness_by_family.png      completeness by family
%  S6  C1_ncrna_call_landscape.png        calls per locus; zero/one/multiple
%  S7  C3_model_composition.png           per-model composition, calls and distinct sequences
%  S8  D3_cds_and_strand.png              intervening CDS and strand detail
%  S9  D5_strand_agreement.png            opposite-strand placements in detail
%  S10 F_pairing_topology.png             partner degrees and component shapes
%  S11 L_dataset_rule_and_multiplicity.png  the stricter cumulative rule (overlaps excluded)
%  S12 N7_burden_and_cooccurrence.png     RT burden per genome and family co-occurrence
%  S13 M3_operons_minimal.png             representative loci, minimal rendering
%
%  NOTE on S11: this is the alternative cumulative rule in which overlapping placements are
%  EXCLUDED. Its step-4 value is 26,293 pairs against 29,440 in Table~\ref{tab:funnel}. The
%  difference is entirely the treatment of overlaps. Do not present the two as versions of
%  one funnel.
%
%  NOTE on S12: co-occurrence is measured on RT-carrying genomes only, so it describes
%  co-occurrence among them and not enrichment. It is an exploratory extension pending
%  assembly-level deduplication.
% =====================================================================
